\section{Limitations}

The current implementation is local by design. The trace inputs are fixtures, the bundles are local directories, and the packages are generated as local research objects. This design supports reproducibility for reviewers because the same local corpus can be validated, resolved, imported, packaged, and evaluated without depending on external services. However, it does not demonstrate operation in a deployed agent platform, live observability stack, or production release pipeline. The reported results should therefore be read as local evidence of conformance behavior, not as deployment evidence.

The implementation also has explicit external integration limits. It does not connect to a real OpenTelemetry collector, call the LangSmith API, submit records to an FDO registry, issue externally resolvable identifiers, or provide full RO-Crate certification. The terms OpenTelemetry-like, LangSmith-like, FDO-like, and RO-Crate-like are used deliberately. They indicate local fixtures, local mappings, and local package structures, not external conformance claims. This boundary avoids overstating the implementation while preserving a clear path for future integration work.

The privacy and security checks are limited. ASIEP uses a reference-and-digest design to avoid embedding raw sensitive artifacts by default, and the current evaluator checks for sentinel strings and sensitive-key patterns. This is useful for detecting common local violations of the reference-only policy, but it is not a substitute for full data-loss prevention, access-control enforcement, cryptographic signature verification, secure key management, or legal compliance review. A production implementation would need stronger redaction, authorization, retention, audit, and incident-response controls.

The evaluation corpus is curated and local. It includes known valid examples, invalid records, tampered bundles, missing evidence cases, digest mismatches, path-escape attempts, unsafe promotion cases, and policy violations. This corpus is sufficient to test whether the current toolchain behaves consistently on known cases, but it is not a large-scale benchmark of agent systems. The reported tamper-detection recall and false-positive rate are measured only on this corpus. They should not be interpreted as evidence of general robustness against all attacks or all forms of evidence manipulation.

ASIEP is also not claimed to be complete for all agent systems. It focuses on self-improvement events and treats them as evidence objects. Other critical agent-to-agent actions, such as delegation, negotiation, data access, transaction execution, cross-organization authorization, or multi-party dispute resolution, may require additional profiles under a broader Agent Accountability Evidence Layer. The current profile does not evaluate human-in-the-loop organizational governance, institutional approval processes, or long-term operational monitoring across real deployments.

Future work should proceed in controlled stages: test ASIEP on live traces, extend the corpus across multiple agent frameworks, add stronger cryptographic verification, evaluate human review workflows, and engage provenance, FDO, RO-Crate, observability, and AI governance communities. These extensions should preserve the same boundary discipline: ASIEP should remain a minimal evidence profile for auditable self-improvement events, not a broad claim to solve agent governance as a whole.
